\chapter{Thyroglobulin}

\textbf{Interpretive caution:} Thyroglobulin is interpreted using the patient's surgical and radioactive-iodine history, TSH level, assay, anti-thyroglobulin antibody status, and imaging. Numeric response categories are not universal reference ranges.

\section*{What Is This Test?}
Thyroglobulin (Tg) is a large 660 kDa dimeric glycoprotein produced exclusively by thyroid follicular cells. It serves as the protein scaffold for synthesis of the thyroid hormones thyroxine (T4) and triiodothyronine (T3). Tyrosine residues within the Tg molecule are iodinated and coupled within the follicular lumen to form T4 and T3, which are then released into the circulation through proteolytic cleavage. Small amounts of intact thyroglobulin leak into the bloodstream under normal conditions. The serum thyroglobulin test measures the concentration of this protein in blood. Its primary clinical application is as a tumor marker for differentiated thyroid carcinoma (DTC) --- papillary (PTC, 80--85\% of all thyroid cancers) and follicular (FTC, 10--15\%) --- after total thyroidectomy and radioactive iodine (I-131) ablation. Under these conditions, any residual or recurrent thyroid tissue (normal or malignant) is the only source of thyroglobulin, so serum Tg becomes a highly sensitive and specific biomarker. Thyroglobulin has no role in screening for thyroid cancer in a general population or in initial diagnosis of thyroid nodules. Its utility is confined to follow-up surveillance after definitive treatment.

\section*{Alternative Names}
Tg; Serum Thyroglobulin; Thyroglobulin Tumor Marker; hTg; Tg by Immunoassay; Thyroglobulin RIA; Thyroglobulin Cancer Marker

\section*{Principle}
Thyroglobulin is measured using immunometric (sandwich) assays on automated platforms (Roche Cobas ECLIA, Beckman Coulter DxI, Siemens IMMULITE, Abbott Architect). Two monoclonal antibodies recognizing different epitopes on the Tg molecule are used: one antibody (capture) is bound to a solid phase and a second antibody (detection) is labeled with a chemiluminescent tag. The assay is calibrated against the certified reference material CRM-457 (BCR/IRMM). Detection limits of high-sensitivity (second-generation) Tg assays are 0.1--0.2 ng/mL. The major limitation is interference from anti-thyroglobulin autoantibodies (TgAb), which are present in approximately 15--25\% of DTC patients. TgAb interference causes falsely low (underestimation) results in immunometric assays because the autoantibodies block epitopes recognized by the assay antibodies. For this reason, Tg measurement must always be accompanied by TgAb measurement on the same specimen. If TgAb is positive, thyroglobulin results are unreliable, and TgAb trends become the surrogate tumor marker --- a decline in TgAb titers suggests disease remission, while rising titers suggest recurrence. Radioimmunoassay (RIA) for Tg is less affected by TgAb interference but has poorer sensitivity and is rarely used. Liquid chromatography-tandem mass spectrometry (LC-MS/MS) is emerging as a TgAb-free method that measures Tg-specific proteolytic peptides, circumventing antibody interference, though it is not yet widely available.

\section*{History}
Thyroglobulin was first purified from thyroid extracts in the 1930s. The first radioimmunoassay for serum thyroglobulin was developed by Andre Van Herle and colleagues in 1973 at UCLA. The utility of Tg as a tumor marker for differentiated thyroid cancer was recognized by Van Herle in 1975. Landmark studies by Spencer and colleagues in the 1980s--1990s demonstrated that Tg is most useful after total thyroid ablation and that TgAb interference is a critical issue in approximately 20\% of patients. The first WHO-certified reference material for Tg (CRM-457) was established in 1994. The development of high-sensitivity immunometric assays with functional sensitivity of 0.1--0.2 ng/mL in the 2000s revolutionized DTC monitoring by enabling Tg detection without the need for stimulated (withdrawal or rhTSH) testing. The 2015 ATA guidelines established the concept of excellent response (Tg <0.2 ng/mL on suppression with negative TgAb and negative imaging) as the lowest risk category, permitting less intensive follow-up. LC-MS/MS methods for Tg were developed in the 2010s to address TgAb interference.

\section*{Why This Test Is Ordered}
\begin{itemize}
  \item Surveillance for persistent or recurrent differentiated thyroid carcinoma (papillary or follicular) after total thyroidectomy with or without radioactive iodine ablation --- the primary indication for thyroglobulin testing
  \item Pre-radioactive iodine ablation assessment in newly diagnosed DTC to establish a post-operative baseline thyroglobulin level (pre-ablation Tg) that predicts risk of persistent disease
  \item Monitoring response to therapy using the ATA risk stratification system: excellent response (Tg <0.2 ng/mL), indeterminate (Tg 0.2--1.0, or higher but declining), biochemical incomplete (Tg rising or persistently elevated without structural disease), and structural incomplete (elevated Tg with positive imaging)
  \item Detection of recurrence in patients previously in remission after DTC treatment --- serial Tg with rising trends is more sensitive than imaging for detecting microscopic recurrence
  \item Evaluation of thyrotoxicosis etiology --- undetectable thyroglobulin in a hyperthyroid patient with suppressed TSH suggests exogenous thyroid hormone ingestion (factitious thyrotoxicosis) since endogenous thyroid production is suppressed
  \item Investigation of ectopic thyroid tissue (lingual thyroid, struma ovarii) --- thyroglobulin may be produced from ectopic tissue; functional struma ovarii can produce Tg, T4, and T3
  \item Assessment of subacute thyroiditis --- markedly elevated Tg (>1000 ng/mL) with suppressed TSH and low RAI uptake supports destructive thyroiditis with follicular leakage of colloid
  \item Neonatal congenital hypothyroidism workup --- absent or very low thyroglobulin in hypothyroid neonates with normal TSH receptor suggests thyroglobulin synthesis defects (dyshormonogenesis); detectable Tg suggests thyroid agenesis (no tissue, no Tg)
  \item TgAb trends as surrogate biomarker in TgAb-positive DTC patients --- declining TgAb predicts remission while rising levels may indicate recurrence
\end{itemize}

\section*{Is This a Primary or Secondary Test}
This is a tumor marker used for surveillance after definitive treatment of differentiated thyroid carcinoma. It is not a primary screening test. Tg has no role in initial thyroid nodule evaluation (this is TSH, ultrasound, and FNA), in screening for thyroid cancer in asymptomatic individuals, or in diagnosing primary thyroid dysfunction (TSH and FT4 serve this role). In DTC follow-up, Tg is the primary biochemical marker and the cornerstone of long-term surveillance.

\section*{How This Test Is Performed}
A healthcare professional draws a blood sample from a vein in your arm into a serum separator tube (gold-top) or plain red-top tube. Approximately 5 mL of blood is collected. The specimen is centrifuged to separate serum. Tg is stable in serum at 2--8 degrees C for 7 days and at -20 degrees C for 6 months. The Tg measurement must always be accompanied by TgAb testing on the same specimen. The test should be performed using the same laboratory and assay platform across serial measurements because there is significant inter-assay variability (5--15\%) and different assays recognize different epitopes, leading to discordant absolute values. When changing assays, a re-baseline measurement should be obtained before the switch.

\section*{What Preparation Is Needed}
No fasting or special preparation is required. Tg is measured on levothyroxine suppression (suppressed Tg) in stable, low-risk patients. For high-risk patients or when recurrence is suspected, stimulated Tg measurement (after thyroid hormone withdrawal to elevate TSH, or after recombinant human TSH [rhTSH, Thyrogen, 0.9 mg intramuscularly for 2 days]) may be ordered. Stimulated testing increases the sensitivity of Tg for detecting small residual disease by increasing TSH-dependent Tg production. rhTSH administration requires coordination: 0.9 mg intramuscular injection on Days 1 and 2, blood sample for Tg on Day 3 (72 hours after first injection). High-dose biotin (>5 mg/day) must be discontinued 48 hours before testing to avoid interference with streptavidin-biotin immunoassay platforms.

\section*{Contraindications and Precautions}
TgAb is an important limitation, not a contraindication to testing. TgAb can interfere with immunometric Tg assays, often causing underestimation, but the magnitude and direction of interference vary by assay; Tg results should therefore be interpreted cautiously rather than automatically discarded. TgAb must be measured with every Tg sample, and serial TgAb trends can provide supportive information when interpreted with imaging. Tg levels are TSH-dependent, so the TSH level at measurement must be known. Patients with a remaining thyroid lobe have endogenous Tg and require a different interpretive framework. Heterophile antibodies and human anti-mouse antibodies can cause falsely elevated results. Because inter-assay variability is clinically important at low concentrations, use the same assay for serial monitoring whenever possible. ATA response categories and Tg thresholds apply only in defined postoperative contexts and should not be treated as universal cutoffs.

\section*{Risks}
Minimal risk from venipuncture (slight pain, bruising, hematoma). Rare: vasovagal syncope or infection. No radiation exposure. For rhTSH-stimulated testing: mild nausea (10\%), headache (8\%), asthenia (5\%), and vomiting (2\%). Allergic reactions to rhTSH are rare.

\section*{What Happens After the Test}
Results are available within 1--4 hours on automated immunoassay platforms. Interpretation follows ATA risk-stratified response categories: (1) Excellent response --- suppressed Tg <0.2 ng/mL, TgAb negative, negative neck ultrasound --- risk of recurrence <1\% over 5--10 years, TSH suppression target may be relaxed to 0.5--2.0 mIU/L, and monitoring frequency reduced to annual; (2) Indeterminate response --- Tg 0.2--1.0 ng/mL or TgAb stable/decreasing --- 15--20\% risk of structural disease requiring additional imaging, continue observation with serial Tg; (3) Biochemical incomplete --- rising Tg or TgAb without structural disease on imaging --- risk of structural recurrence 20--40\%, additional imaging (neck ultrasound, CT chest, PET-CT if Tg >10 ng/mL) needed; (4) Structural incomplete --- persistent or new disease on imaging regardless of Tg level --- risk of additional recurrence >50\%, further treatment needed (surgery, RAI, external beam radiation, tyrosine kinase inhibitors if RAI-refractory).

\section*{Understanding Results}

\subsection*{Below Normal}
\begin{itemize}
  \item \textbf{Undetectable Tg (<0.2 ng/mL) on suppression with negative TgAb:} Excellent response after total thyroidectomy and RAI ablation for DTC. Indicates complete remission with <1\% risk of recurrence at 5--10 years. TSH suppression can be relaxed to 0.5--2.0 mIU/L.
  \item \textbf{Undetectable Tg with positive TgAb:} Cannot be interpreted as remission because TgAb may mask low-level Tg. TgAb trends become the biomarker: declining TgAb over 1--3 years with negative imaging is considered a surrogate for remission; stable or rising TgAb warrants imaging.
  \item \textbf{Undetectable Tg in congenital hypothyroidism:} May support thyroid agenesis or athyreosis, but must be interpreted with thyroid imaging, TSH, and the clinical evaluation; a detectable result suggests the presence of thyroid tissue but does not by itself establish the cause of congenital hypothyroidism.
  \item \textbf{Very low Tg (<1 ng/mL) in hyperthyroid patient:} Suggests exogenous thyrotoxicosis (factitious or therapeutic levothyroxine ingestion) --- endogenous thyroid function is suppressed, so little or no Tg is produced. This contrasts with all forms of endogenous hyperthyroidism (Graves, toxic nodular goiter, subacute thyroiditis) where Tg is elevated.
\end{itemize}

\subsection*{Above Normal}
\begin{itemize}
  \item \textbf{Detectable or rising Tg after total thyroidectomy and RAI:} Strongly suggests residual thyroid tissue or recurrent/metastatic differentiated thyroid carcinoma. Tg levels correlate with tumor burden: Tg 1--10 ng/mL suggests minimal residual disease; Tg 10--100 ng/mL suggests moderate tumor burden; Tg >100 ng/mL suggests significant metastatic disease (often pulmonary or skeletal). Stimulated Tg >2 ng/mL has approximately 80--90\% sensitivity for detecting structural disease.
  \item \textbf{Even minimally detectable Tg (0.2--1.0 ng/mL) on suppression:} Classified as indeterminate response. Approximately 15--20\% of these patients will have structural disease on further imaging. Rising trends (doubling time <1 year) are more concerning than stable low levels. Doubling time >3 years is a favorable prognostic indicator.
  \item \textbf{Elevated Tg in intact thyroid (no thyroidectomy):} Nonspecific --- occurs in goiter, thyroiditis, Graves disease, toxic nodular goiter, thyroid adenomas, and thyroid cancer. Tg is NOT used as a screening test in patients with an intact thyroid gland. Mildly elevated Tg (20--100 ng/mL) is common in benign nodular goiter; markedly elevated Tg (>1000 ng/mL) is seen in destructive thyroiditis or massive metastatic DTC.
  \item \textbf{Elevated Tg in subacute thyroiditis:} Destructive release of follicular colloid contents including preformed Tg, T4, and T3. Tg levels are often >1000 ng/mL and parallel the rise and fall of T4/T3. TSH is suppressed, and RAI uptake is very low.
  \item \textbf{Elevated Tg in Graves disease:} Increased thyroid parenchymal mass and TSH receptor stimulation drive excess Tg production and release.
  \item \textbf{Elevated Tg in endemic goiter:} Iodine deficiency drives TSH elevation and thyroid hyperplasia, resulting in increased Tg production correlating with goiter size.
  \item \textbf{Elevated Tg in congenital dyshormonogenesis:} Defects in Tg synthesis, iodination, or coupling lead to elevated TSH and compensatory Tg production. Detectable Tg with hypothyroidism and goiter suggests dyshormonogenesis rather than agenesis (where Tg is absent).
  \item \textbf{Rising TgAb in DTC patient with initially positive TgAb:} Considered biochemical recurrence even if unstimulated Tg is undetectable. A rising TgAb titer over time, confirmed by at least 2 measurements, merits further imaging investigation.
\end{itemize}

\section*{Normal Results}
\begin{itemize}
  \item \textbf{Post-total thyroidectomy with RAI ablation (on suppression):} A value below the assay's functional sensitivity, often around 0.1--0.2 ng/mL, may support an excellent response when TgAb is negative and imaging is negative; this is not a universal cutoff.
  \item \textbf{Post-total thyroidectomy with RAI ablation (rhTSH-stimulated):} <1.0 ng/mL (undetectable or <1.0) --- stimulated Tg <1 ng/mL predicts <1\% risk of structural disease at 5 years
  \item \textbf{Intact thyroid gland (healthy euthyroid):} 1.5--30 ng/mL (varies with iodine intake and thyroid volume)
  \item \textbf{Intact thyroid (iodine-sufficient regions):} 3--40 ng/mL (upper limit higher with iodine deficiency)
  \item \textbf{TgAb status:} Use the laboratory's assay-specific reference interval. Positive TgAb can make immunometric Tg results unreliable; serial TgAb trends may provide additional information but are not a standalone diagnosis of recurrence.
  \item \textbf{Pre-ablation Tg (after thyroidectomy, before RAI):} Variable --- higher levels (>5--10 ng/mL) predict greater risk of residual disease
  \item \textbf{Note:} Thyroglobulin reference intervals apply only to patients with intact thyroid glands. After thyroidectomy, Tg must be interpreted in relation to radioactive iodine treatment, TSH stimulation, TgAb, imaging, and the patient's risk category; detectability alone does not establish recurrence. Use the same assay and laboratory for serial measurements whenever possible.
\end{itemize}

\section*{Follow-up Tests}
\begin{itemize}
  \item \textbf{Anti-thyroglobulin antibodies (TgAb):} Measure TgAb with every Tg measurement. TgAb can interfere with immunometric Tg results, so a Tg result may be unreliable when TgAb is present. Serial TgAb trends, interpreted with imaging and the clinical context, can provide supportive information.
  \item \textbf{Neck ultrasound (high-resolution, 7--15 MHz):} First-line imaging for DTC recurrence. Evaluates thyroid bed for recurrence and cervical lymph node compartments (levels I--VI) for metastatic lymphadenopathy. Abnormal lymph nodes (loss of fatty hilum, microcalcifications, cystic changes, peripheral vascularity) prompt fine-needle aspiration with Tg washout.
  \item \textbf{Diagnostic whole-body radioactive iodine scan (WBS, I-123 or I-131):} Only if patient has proven iodine-avid disease (previous positive post-therapy scan) and negative Tg with rising Tg or concerning ultrasound. Not routinely used if stimulated Tg is <1 ng/mL and ultrasound is negative.
  \item \textbf{Neck and chest CT with contrast:} For patients with elevated Tg and negative neck ultrasound, to detect mediastinal lymph node metastases or pulmonary micrometastases (miliary pattern on CT).
  \item \textbf{FDG-PET/CT:} Indicated for Tg >10 ng/mL with negative I-131 whole-body scan (TENIS syndrome --- Tg-Elevated Negative Iodine Scan), suggesting dedifferentiated/R AI-refractory disease. FDG-avid disease carries a worse prognosis.
  \item \textbf{Fine-needle aspiration (FNA) of suspicious lymph nodes with Tg washout:} Tg concentration in needle washout fluid >10 ng/mL is highly specific for metastatic lymph node involvement.
  \item \textbf{Recombinant human TSH (rhTSH)-stimulated Tg:} In selected high-risk patients or when unstimulated Tg is borderline, stimulated Tg increases sensitivity for detecting minimal residual disease.
  \item \textbf{Staging CT chest without contrast and bone scan:} For patients with elevated post-operative Tg (>10--30 ng/mL) to evaluate for distant metastases (lungs, bones, liver).
  \item \textbf{Thyroglobulin LC-MS/MS:} Emerging technology that may overcome TgAb interference in DTC follow-up. Not yet widely available.
\end{itemize}

\section*{References}
\begin{enumerate}
  \item Haugen BR, Alexander EK, Bible KC, et al. 2015 American Thyroid Association management guidelines for adult patients with thyroid nodules and differentiated thyroid cancer. Thyroid. 2016;26(1):1-133.
  \item Spencer CA, LoPresti JS. Measuring thyroglobulin and thyroglobulin autoantibody in patients with differentiated thyroid cancer. Nat Clin Pract Endocrinol Metab. 2008;4(4):223-233.
  \item Tuttle RM, Alzahrani AS. Risk stratification in differentiated thyroid cancer: from detection to final follow-up. J Clin Endocrinol Metab. 2019;104(9):4087-4100.
  \item Giovanella L, Clark PM, Chiovato L, et al. Thyroglobulin measurement using highly sensitive assays in patients with differentiated thyroid cancer: a clinical position paper. Eur J Endocrinol. 2014;171(2):R33-46.
  \item Spencer CA, Bergoglio LM, Kazaros E, et al. Serum thyroglobulin autoantibodies: prevalence, influence on serum thyroglobulin measurement, and prognostic significance in patients with differentiated thyroid carcinoma. J Clin Endocrinol Metab. 1998;83(4):1121-1127.
  \item Van Herle AJ, Uller RP, Matthews NL, Brown J. Radioimmunoassay for measurement of thyroglobulin in human serum. J Clin Invest. 1973;52(6):1320-1327.
  \item Baloch Z, Carayon P, Conte-Devolx B, et al. Laboratory medicine practice guidelines: laboratory support for the diagnosis and monitoring of thyroid disease. Thyroid. 2003;13(1):3-126.
  \item Burtis CA, Ashwood ER, Bruns DE. Tietz Textbook of Clinical Chemistry and Molecular Diagnostics. 6th ed. Elsevier; 2023.
\end{enumerate}

\clearpage
